\section{\texorpdfstring{$K$}{K}-微分的函子性}

\begin{proposition}{$K$-微分的自然性}{natuality-of-K-differential}
	设如下关于交换环的图表交换,$\eta,\eta^{\prime}$分别为$A,A^{\prime}$配备$K$-代数,$K^{\prime}$-代数的结构.
	\begin{center}
		\begin{tikzcd}
			A & {A^{\prime}} \\
			K & {K^{\prime}}
			\arrow["\rho", from=1-1, to=1-2]
			\arrow["\eta"', from=1-1, to=2-1]
			\arrow["{\eta^{\prime}}", from=1-2, to=2-2]
			\arrow["u"', from=2-1, to=2-2]
		\end{tikzcd}
	\end{center}
	则存在唯一的$A$-线性映射$v:\Omega_{K}\left(A\right)\to u_{\ast}\left(\Omega_{K^{\prime}}\left(A^{\prime}\right)\right)$使下图交换.
	\begin{center}
		\begin{tikzcd}
			A & {u_{\ast}\left(A^{\prime}\right)} \\
			{\Omega_{K}\left(A\right)} & {u_{\ast}\left(\Omega_{K^{\prime}}\left(A^{\prime}\right)\right)}
			\arrow["u", from=1-1, to=1-2]
			\arrow["{d_{A/K}}"', from=1-1, to=2-1]
			\arrow["{d_{A^{\prime}/K^{\prime}}}", from=1-2, to=2-2]
			\arrow["{\exists!v}"', from=2-1, to=2-2]
		\end{tikzcd}
	\end{center}
\end{proposition}

\begin{definition}{}{}
	沿用命题(\cref{prop:natuality-of-K-differential})的设定.我们把$v$记作$\Omega\left(u\right)$.
\end{definition}

\begin{proposition}{$K$-微分的函子性}{}
	设如下关于交换环的图表交换,$\eta,\eta^{\prime},\eta^{\prime\prime}$分别为$A,A^{\prime},A^{\prime\prime}$配备$K$-代数,$K^{\prime}$-代数,$K^{\prime\prime}$-代数的结构.
	\begin{center}
		\begin{tikzcd}
			A & {A^{\prime}} & {A^{\prime\prime}} \\
			K & {K^{\prime}} & {K^{\prime\prime}}
			\arrow["\rho", from=1-1, to=1-2]
			\arrow["\eta"', from=1-1, to=2-1]
			\arrow["\sigma", from=1-2, to=1-3]
			\arrow["{\eta^{\prime}}"', from=1-2, to=2-2]
			\arrow["{\eta^{\prime\prime}}", from=1-3, to=2-3]
			\arrow["u"', from=2-1, to=2-2]
			\arrow["v"', from=2-2, to=2-3]
		\end{tikzcd}
	\end{center}
	那么,如下图表交换.
	\begin{center}
		\begin{tikzcd}
			A & {\rho_{\ast}\left(A^{\prime}\right)} & {\sigma_{\ast}\left(\rho_{\ast}\left(A^{\prime\prime}\right)\right)} \\
			{\Omega_{K}\left(A\right)} & {u_{\ast}\left(\Omega_{K^{\prime}}\left(A^{\prime}\right)\right)} & {v_{\ast}\left(u_{\ast}\left(\Omega_{K^{\prime\prime}}\left(A^{\prime\prime}\right)\right)\right)}
			\arrow["\rho", from=1-1, to=1-2]
			\arrow["\eta"', from=1-1, to=2-1]
			\arrow["\sigma", from=1-2, to=1-3]
			\arrow["{\eta^{\prime}}"', from=1-2, to=2-2]
			\arrow["{\eta^{\prime\prime}}", from=1-3, to=2-3]
			\arrow["{\Omega\left(u\right)}"', from=2-1, to=2-2]
			\arrow["{\Omega\left(v\circ u\right)}", curve={height=30pt}, from=2-1, to=2-3]
			\arrow["{\Omega\left(v\right)}"', from=2-2, to=2-3]
		\end{tikzcd}
	\end{center}
\end{proposition}

\begin{definition}{}{}
	沿用命题(\cref{prop:natuality-of-K-differential})的设定.记$\Omega_{0}\left(u\right):u^{\ast}\left(\Omega_{K}\left(A\right)\right)\to\Omega_{K^{\prime}}\left(A^{\prime}\right)$是$u$诱导的$A^{\prime}$-线性映射.
\end{definition}

\begin{proposition}{基变换下的$K$-微分模同构}{}
	沿用命题(\cref{prop:natuality-of-K-differential})的设定.设$A^{\prime}=u^{\ast}\left(K^{\prime}\right)$,$\eta^{\prime}$和$u$是典范同态,则$\Omega_{0}\left(u\right)$是$A^{\prime}$-模同构.
\end{proposition}

接下来的内容中,默认$K=K^{\prime},\rho=\id_{K},B=A^{\prime},u\in\Hom_{K-\mathsf{Alg}}\left(A,B\right)$.

\begin{proposition}{$K$-微分模的第一基本正合列}{}
	设$\Omega_{u}:\Omega_{K}\left(B\right)\to\Omega_{A}\left(B\right)$是典范$B$-线性映射,则如下序列正合(称之为第一基本正合列).
	\begin{align*}
		u_{\ast}\left(u^{\ast}\left(\Omega_{K}\left(A\right)\right)\right)\xlongrightarrow{\Omega_{0}\left(u\right)}u_{\ast}\left(\Omega_{K}\left(B\right)\right)\xlongrightarrow{\Omega_{u}}u_{\ast}\left(\Omega_{A}\left(B\right)\right)\xlongrightarrow{}0
	\end{align*}
\end{proposition}

\begin{proposition}{$K$-微分模的第二基本正合列}{}
	设$\mathfrak{J}$是$A$的理想,$B=A/\mathfrak{J}$,$u$是满射.考虑典范嵌入$i_{A}:\Omega_{K}\left(A\right)\to u_{\ast}\left(u^{\ast}\left(\Omega_{K}\left(A\right)\right)\right)$和$d^{\prime}\coloneq i_{A}\circ d|_{\mathfrak{J}}$诱导的$B$-线性映射
	\begin{align*}
		\bar{d}:\mathfrak{J}/\mathfrak{J}^{2}\to u^{\ast}\left(\Omega_{K}\left(A\right)\right),
	\end{align*}
	则$B$-线性映射的序列$\mathfrak{J}/\mathfrak{J}^{2}\xlongrightarrow{\bar{d}}u^{\ast}\left(\Omega_{K}\left(A\right)\right)\xlongrightarrow{\Omega_{0}\left(u\right)}\Omega_{K}\left(B\right)\xlongrightarrow{}0$正合(称为第一正合列).
\end{proposition}

\begin{proposition}{局部化下的$K$-微分模同构}{}
	设$A$作为环是整环,$B=\Frac\left(A\right)$,$u$是典范映射,则$\Omega_{0}\left(u\right)$是$B$-模同构.
\end{proposition}


























